\section{Методологія розробки}
\label{sec:metodolohiia}

\subsection{Розробка на основі специфікації (SDD)}

Система розробляється відповідно до методології Specification Driven Development (SDD). Специфікація є контролюючим джерелом поведінки системи~--- жодний вихідний код не може визначати поведінку незалежно від цього документа. \sked{K53}

Обов'язковий процес при кожній зміні поведінки системи:

\begin{enumerate}
  \item Оновити цей документ до внесення змін до коду.
  \item Призначити або оновити ідентифікатор вимоги (REQ-F-NNN або REQ-NF-NNN).
  \item Визначити критерій прийняття (AC-NNN) для кожної вимоги.
  \item Написати автоматизований тест до реалізації (TDD, \secref{sec:metodolohiia}).
  \item Запустити тест і переконатися, що він провалюється з очікуваної причини (Red).
  \item Реалізувати мінімальний код, що проходить тест (Green).
  \item Рефакторити лише після успішного проходження тестів (Refactor).
  \item Оновити log-артефакти після завершення ітерації.
\end{enumerate}

\subsection{Обов'язкова розробка через тестування (TDD)}

TDD є обов'язковою для усієї реалізаційної роботи. Цикл Red~/ Green~/ Refactor застосовується до кожної функціональної зміни: \sked{K53}

\begin{itemize}
  \item \textbf{Red}~--- спочатку написати тест, що провалюється.
  \item \textbf{Green}~--- реалізувати мінімальний код, що проходить тест.
  \item \textbf{Refactor}~--- покращити код без зміни зовнішньої поведінки.
\end{itemize}

Жоден виробничий код не приймається, якщо він не покритий принаймні одним автоматизованим тестом. Реальні зовнішні системи (СУБД, сторонні API) не використовуються в unit-тестах~--- застосовуються детерміновані замінники (stubs / mocks).

\subsection{Правило трасованості}

Кожна реалізована поведінка системи повинна мати повну трасованість у ланцюжку:

\begin{codelisting}
\caption{Ланцюжок трасованості}
\label{lst:trace}
\begin{skedcode}
REQ → AC → Test → Implementation
\end{skedcode}
\end{codelisting}

Функціональність без ідентифікатора вимоги не є частиною системи. Вимога без критерію прийняття не готова до реалізації.

\subsection{Log-артефакти тестування}

Після кожного запуску тестів агент-розробник записує log-файл із результатами. \sked{K54} Log-файл повинен містити:

\begin{itemize}
  \item загальну кількість тестів і кількість passed~/ failed;
  \item статус кожного окремого тесту (passed~/ failed~/ skipped);
  \item час виконання тестового набору.
\end{itemize}

Log-файл є машиночитаним артефактом верифікації~--- на його основі агент-розробник і людина-рецензент визначають, чи виконане завдання відповідно до специфікації (крок «Is~verified?» у SDLC-процесі).
